\section{RESULTS AND ANALYSIS}

Table~\ref{tab:main-results} reports the complete three-seed NER evidence. The
domain block is the primary frozen-test evaluation; CLUENER tests directional
recurrence with separately constructed resources.

\begin{table}[ht]
\caption{NER results at $K=5$ over seeds 42, 43, and 44. Domain uses the
frozen test set; CLUENER uses all labeled official development rows. Values are
mean $\pm$ sample standard deviation.}
\label{tab:main-results}
\centering
\scriptsize
\setlength{\tabcolsep}{2.1pt}
\begin{tabular}{ccccc}
\toprule
Config. & Precision & Recall & F1 & FP count \\
\midrule
\multicolumn{5}{l}{\textit{Equipment-domain frozen test}} \\
A & $0.484066\pm0.025428$ & $0.599231\pm0.028904$ & $0.534645\pm0.004375$ & $613.3\pm88.8$ \\
B & $0.505743\pm0.000234$ & $0.615304\pm0.001816$ & $0.555169\pm0.000880$ & $573.7\pm1.2$ \\
C & $0.511433\pm0.000472$ & $0.640811\pm0.001210$ & $0.568858\pm0.000769$ & $584.0\pm0.0$ \\
D & $0.525135\pm0.000000$ & $0.613208\pm0.000000$ & $0.565764\pm0.000000$ & $529.0\pm0.0$ \\
\midrule
\multicolumn{5}{l}{\textit{CLUENER labeled development}} \\
A & $0.448636\pm0.036890$ & $0.413086\pm0.056862$ & $0.428519\pm0.037917$ & $1569.0\pm306.8$ \\
B & $0.528488\pm0.001929$ & $0.500217\pm0.002118$ & $0.513964\pm0.001907$ & $1371.0\pm7.2$ \\
C & $0.530932\pm0.001140$ & $0.504774\pm0.001316$ & $0.517521\pm0.000619$ & $1370.0\pm8.7$ \\
D & $0.540764\pm0.002162$ & $0.488715\pm0.000940$ & $0.513424\pm0.001493$ & $1275.0\pm8.7$ \\
\bottomrule
\end{tabular}
\end{table}

\subsection{Equipment-Domain Operating Points}

kNN demonstrations improve mean F1 by 0.020524 from A to B and reduce FP by
approximately 39.7. Adding the glossary raises F1 by another 0.013689; recall
increases from 0.615304 to 0.640811, while FP rises by about 10.3. The glossary
therefore acts primarily as a coverage and recall control.

C and D are near-tied in F1 (0.568858 versus 0.565764). Relative to C, D raises
precision from 0.511433 to 0.525135 and reduces mean FP from 584.0 to 529.0,
while recall falls from 0.640811 to 0.613208. Thus, RAG changes the operating
point rather than maximizing F1: C favors recall, whereas D favors precision
and a smaller review burden. These differences are descriptive.

\subsection{Focused CLUENER Replication}

On CLUENER, B improves F1 over A by 0.085445. The support-derived lexicon adds
a modest 0.003557 from B to C, accompanied by higher recall. D remains near C
in F1 ($-0.004097$), increases precision by 0.009832, and removes 95 FP while
reducing recall by 0.016059. The same module directions therefore recur on the
complete designated labeled-dev block. Because resources are schema-specific
and the official test labels are unavailable, this result is limited
descriptive replication rather than broad cross-domain evidence.

\subsection{False-Positive Analysis}

The locked seed-42 comparison shows that D avoids 181 C-only FP and introduces
126 D-only FP, for the observed net reduction of 55. Table~\ref{tab:fp-categories}
groups the avoided errors; Supplementary Table S4 provides representative
audited cases.

\begin{table}[ht]
\caption{Seed-42 categories of C-only false positives avoided by D.}
\label{tab:fp-categories}
\centering
\begin{tabular}{lr}
\toprule
Error category & Count \\
\midrule
Unsupported entity & 84 \\
Type or nested-overlap error & 34 \\
Over-labeled short alias or code & 25 \\
Boundary error & 24 \\
Metric fragment & 14 \\
\bottomrule
\end{tabular}
\end{table}

RAG therefore acts as an imperfect constraint, not an error eliminator. The
audit supports the observed net FP reduction and higher precision in this
setting; it does not convert FP count into a general factuality measure.

\subsection{Exploratory Relation Boundary}

The separate seed-42 Stage-R audit yields exact relation F1 of
0.040541/0.056972/0.057971/0.045381 for A/B/C/D. D is below C, and every cell
has zero \texttt{HAS\_METRIC} true positives. These results support the
traceability of the two-stage route, but not positive RE or RAG-for-RE
effectiveness. Supplementary Tables S5--S6 retain the full relation counts,
parse outcomes, and development diagnostics. The bounded verifier-deletion
audits also remained below their prespecified positive-evidence threshold, so
verification is not part of the supported final-test contribution.
